% !TEX root = ../main.tex
\chapter{信号から計算へ}
\label{chap:digital-computer}

\begin{chapterabstract}
送信ボタンは押されましたが、要求はまだ端末を出ていません。指やキーボードの操作を受けた回路では電圧が変化し、その変化が入力装置からCPUへ伝わろうとしています。ここで視点を最も小さな世界まで近づけます。

連続的な電圧から安定した0と1を取り出し、トランジスタを論理ゲートへ、論理ゲートを加算器と記憶回路へ組み上げます。さらに、限られた速い記憶と大きな遅い記憶をつなぎ、CPUが一つの命令を実行するところまで進みます。章の終わりには、画面上の操作も無数の状態変化として動いていることが見えるはずです。
\end{chapterabstract}


\section{電圧と論理値}
\label{sec:01-digital-computer-01}

送信の瞬間へカメラを寄せると、最初に見えるのは文字でも命令でもなく、導線や素子に現れる電圧です。電圧は連続的に揺れますが、そのままでは安定した計算の部品にできません。そこで、ある範囲をLow、別の範囲をHighと読むところから始めます。

\begingroup
\begingroup
\paragraph{デジタル抽象化}
        回路は連続量である電圧を、LowとHighの二つの状態として解釈します。
        「Lowは必ず0 V、Highは必ず5 V」という固定値ではありません。

\paragraph{しきい値}
        入力電圧が $V_{IL}$ 以下ならLow、$V_{IH}$ 以上ならHighと判定します。
        二つのしきい値の間は、論理値を保証しない領域です。

\par
\endgroup
\medskip
\begingroup
      \centering
      \begin{tikzpicture}[x=0.75mm,y=0.75mm]
        \draw[->,draw=DeepNavy] (0,0) -- (65,0) node[right] {電圧};
        \fill[PaleBlue] (0,5) rectangle (20,20);
        \fill[RuleGray] (20,5) rectangle (43,20);
        \fill[PaleBlue] (43,5) rectangle (65,20);
        \node at (10,12.5) {Low};
        \node[font=\scriptsize] at (31.5,12.5) {保証外};
        \node at (54,12.5) {High};
        \draw[draw=MutedBlue] (20,3) -- (20,22) node[above] {$V_{IL}$};
        \draw[draw=MutedBlue] (43,3) -- (43,22) node[above] {$V_{IH}$};
      \end{tikzpicture}
      \vspace{4mm}
\paragraph{ノイズ耐性}
        出力がしきい値から十分離れるように設計すると、小さな電圧変動が論理値の反転につながりにくくなります。

\par
\endgroup
  \source{MIT 6.004, “The Digital Abstraction”}
\par
\endgroup


\section{2進数}
\label{sec:01-digital-computer-02}

LowとHighを区別できても、一つの状態が表せるのは二通りだけです。入力された文字や注文番号を扱うには、二状態を順序よく並べ、桁ごとに重みを与える必要があります。ここで電気的な二状態は、数を記すための0と1になります。

\begingroup
\begingroup
\paragraph{ビット、バイト、ワード}
        \term{ビット}は0または1の一桁で、通常は8ビットをまとめて\term{バイト}として扱います。
        \term{ワード}はCPUが基本単位として扱うビット幅を指し、その幅は命令セットによって異なります。

\paragraph{同じ値を異なる基数で表す例}
        10進数の45は、2進数と16進数では次のように表せます。
        \[
          45_{10}=0010\,1101_2=2\mathrm{D}_{16}
        \]
        16進数一桁は4ビットに対応するため、ビット列を短く記述できます。

\par
\endgroup
\medskip
\begingroup
\paragraph{桁の重み}
        6ビットの各桁へ重みを対応づけると、値45を次のように確認できます。
        \centering
        \begin{tabular}{c@{\quad}c@{\quad}c@{\quad}c@{\quad}c@{\quad}c@{\quad}c}
          桁 & 5 & 4 & 3 & 2 & 1 & 0\\
          重み & 32 & 16 & 8 & 4 & 2 & 1\\
          値 & 1 & 0 & 1 & 1 & 0 & 1\\
          \bottomrule
          \multicolumn{7}{c}{$32+8+4+1=45$}
        \end{tabular}

\paragraph{変換の確認}
        10進数から2進数へは2で割った余りを逆順に並べます。
        2進数から10進数へは、1になっている桁の重みを加えます。

\par
\endgroup
\par
\endgroup


\section{負数と2の補数}
\label{sec:01-digital-computer-03}

桁を並べれば非負の数は表せますが、実際の計算には差や増減も現れます。減算専用のまったく別の回路を増やす代わりに、同じ加算回路で負数まで扱える表現が2の補数です。

\begingroup
\begingroup
\paragraph{符号なし整数}
        $n$ビットの範囲は $0$ から $2^n-1$ です。
        8ビットなら $0$ から $255$ までを表します。

\paragraph{2の補数}
        $n$ビットで $-x$ を表すビット列は、$2^n-x$ の下位$n$ビットです。
        範囲は $-2^{n-1}$ から $2^{n-1}-1$ です。

\par
\endgroup
\medskip
\begingroup
\paragraph{8ビットで $-5$ を作る}
        2の補数表現は、正の値をビット反転して1を加えると確認できます。
        \begin{enumerate}
          \item $5=0000\,0101_2$ と表します。
          \item 各ビットを反転して $1111\,1010_2$ にします。
          \item 1を加えて $1111\,1011_2$ を得ます。
        \end{enumerate}
        $0000\,0101_2+1111\,1011_2$ の下位8ビットは0です。

\paragraph{オーバーフロー}
        固定幅を超えた桁は表現できません。
        符号なし加算の桁上がりと、符号付き加算のオーバーフローは別の条件です。

\par
\endgroup
\par
\endgroup


\section{ビット列の解釈}
\label{sec:01-digital-computer-04}

ここまでの0と1には、まだ「数」という意味さえ固定されていません。同じ並びを注文番号、入力文の一文字、小数、命令のどれとして読むかは、ビット列の外側にある規則が決めます。物理的な状態とアプリケーション上の意味をつなぐ最初の境界です。

\begingroup
\begingroup
\paragraph{文字コード}
        Unicodeは文字にコードポイントを割り当てます。
        UTF-8は、そのコードポイントを1〜4バイトの列へ符号化する方式です。

\paragraph{区別する三つの値}
        文字「A」、コードポイント \code{U+0041}、UTF-8のバイト \code{0x41} は、この例では対応します。
        ただし、すべての文字が1バイトになるわけではありません。

\par
\endgroup
\medskip
\begingroup
\paragraph{浮動小数点数}
        代表的な形式は、符号、指数部、仮数部にビットを分けます。
        限られたビット数で広い範囲を表せる一方、表現できる値は離散的です。

\paragraph{丸め誤差}
        10進小数の0.1は、2進数では有限桁に収まりません。
        そのため、\code{0.1 + 0.2} を厳密な10進数の0.3として比較すると失敗する処理系があります。

\par
\endgroup
  \takeaway{ビット列だけでは意味が決まりません。型と符号化方式が解釈を決めます。}
\par
\endgroup


\section{トランジスタ}
\label{sec:01-digital-computer-05}

表現の規則が決まっても、それだけでは状態は変化しません。電圧を受け取り、別の電圧を制御する小さな素子が必要です。抽象的な0と1を現実の動作へ戻すと、そこでトランジスタに出会います。

\begingroup
\begingroup
\paragraph{スイッチとしての抽象化}
        MOSトランジスタでは、ゲートへ加える電圧が端子間の導通を制御します。
        デジタル回路では、この連続的な特性をオン／オフのスイッチとして抽象化します。

\paragraph{CMOS回路}
        特性の異なるnMOSとpMOSを組み合わせると、入力に応じて出力をLowまたはHighへ駆動できます。

\par
\endgroup
\medskip
\begingroup
      \centering
      \begin{tikzpicture}[node distance=4mm,every node/.append style={font=\small}]
        \node[box,minimum width=22mm,minimum height=15mm] (sw) {制御される\\スイッチ};
        \node[box,left=of sw] (in) {入力電圧};
        \node[box,right=of sw] (out) {出力電圧};
        \node[accentbox,above=of sw] (gate) {ゲート電圧};
        \draw[flow] (in) -- (sw);
        \draw[flow] (sw) -- (out);
        \draw[flow] (gate) -- (sw);
      \end{tikzpicture}
      \vspace{1mm}
\paragraph{回路設計で残る現実}
        遅延、消費電力、発熱、しきい値、配線容量は消えません。
        論理回路は、これらの物理特性が仕様範囲内にあることを前提にします。

\par
\endgroup
  \source{MIT 6.004, computation structures lecture notes}
\par
\endgroup


\section{論理ゲート}
\label{sec:01-digital-computer-06}

一個のトランジスタは単純なスイッチとして振る舞います。複数を組み合わせると、「両方がHighならHigh」「どちらかがHighならHigh」という判断を回路自身に行わせられます。このまとまりが、計算を組み立てる語彙になります。

\begingroup
\begingroup
\paragraph{基本ゲートの真理値表}
        AND、OR、XORは、二つの入力に対する出力を次の真理値表で定義できます。
        \centering
        \begin{tabular}{ccccc}
          \toprule
          $A$ & $B$ & NOT $A$ & $A$ AND $B$ & $A$ OR $B$\\
          \midrule
          0 & 0 & 1 & 0 & 0\\
          0 & 1 & 1 & 0 & 1\\
          1 & 0 & 0 & 0 & 1\\
          1 & 1 & 0 & 1 & 1\\
          \bottomrule
        \end{tabular}

\paragraph{組み合わせ回路}
        現在の出力が現在の入力だけで決まる回路です。
        同じ入力には同じ出力が対応します。

\par
\endgroup
\medskip
\begingroup
\paragraph{NAND}
        NANDはANDの結果を反転します。
        \[
          A\mathbin{\mathrm{NAND}}B=\overline{A\land B}
        \]

\paragraph{真理値表の役割}
        回路記号の見た目ではなく、すべての入力組合せに対する出力を仕様として固定します。

\par
\endgroup
\par
\endgroup


\section{NANDによる回路構成}
\label{sec:01-digital-computer-07}

論理ゲートの種類を増やすほど、できることは増えるように見えます。しかし重要なのは種類の多さではなく、必要な論理を組み立てられることです。NAND一種類からほかの論理を作る過程を見ると、単純な部品の組み合わせが複雑さを生む原理が分かります。

\begingroup
\begingroup
\paragraph{NOTを作る}
        同じ信号を二つの入力へ入れます。
        \[
          \overline{A}=A\mathbin{\mathrm{NAND}}A
        \]

\paragraph{ANDを作る}
        NANDの出力を、NANDで再び反転します。
        \[
          A\land B=\overline{\overline{A\land B}}
        \]

\par
\endgroup
\medskip
\begingroup
\paragraph{ORを作る}
        De Morganの法則を使います。
        \[
          A\lor B=\overline{\overline{A}\land\overline{B}}
        \]

\paragraph{万能ゲートという意味}
        NOT、AND、ORで任意のブール関数を表せます。
        それらをNANDで置き換えられるため、NANDだけでも任意の組み合わせ論理を構成できます。

\par
\endgroup
\par
\endgroup


\section{加算器}
\label{sec:01-digital-computer-08}

判断の部品がそろったので、初めて具体的な計算へ進めます。二つのビットを足すには、和だけでなく次の桁へ渡す桁上がりも必要です。小さな論理の出力が次の論理の入力になる様子は、この先のすべての層に繰り返し現れます。

\begingroup
\begingroup
\paragraph{半加算器}
        1ビットの$A$と$B$を加えます。
        \[
          S=A\oplus B,\qquad C=A\land B
        \]
        $S$は和、$C$は上位桁への桁上がりです。

\paragraph{不足する入力}
        半加算器は、下位桁から届く桁上がり$C_{in}$を扱えません。

\par
\endgroup
\medskip
\begingroup
\paragraph{全加算器}
        入力$a$、$b$と下位桁からの桁上がり$c_{in}$に対し、和と次の桁上がりは次式で求められます。
        \[
          S=A\oplus B\oplus C_{in}
        \]
        \[
          C_{out}=(A\land B)\lor(C_{in}\land(A\oplus B))
        \]

      \centering
      \begin{tikzpicture}[node distance=7mm]
        \node[accentbox] (fa) {全加算器};
        \node[box,left=of fa] (inputs) {$A,B,C_{in}$};
        \node[box,right=of fa] (outputs) {$S,C_{out}$};
        \draw[flow] (inputs) -- (fa);
        \draw[flow] (fa) -- (outputs);
      \end{tikzpicture}
\par
\endgroup
\par
\endgroup


\section{多ビット加算とALU}
\label{sec:01-digital-computer-09}

一桁の加算器を横へつなぐと、注文番号やメモリアドレスを扱える幅の計算になります。さらに演算を選ぶ回路を加えると、加算、比較、論理演算を一か所で担うALUになります。CPUの「計算する部分」が姿を見せ始めます。

\begingroup
  \centering
  \begin{tikzpicture}[node distance=5mm]
    \node[box] (f0) {FA 0};
    \node[box,right=of f0] (f1) {FA 1};
    \node[box,right=of f1] (f2) {FA 2};
    \node[box,right=of f2] (f3) {FA 3};
    \draw[flow] (f0) -- node[above,font=\scriptsize] {$C_1$} (f1);
    \draw[flow] (f1) -- node[above,font=\scriptsize] {$C_2$} (f2);
    \draw[flow] (f2) -- node[above,font=\scriptsize] {$C_3$} (f3);
    \node[below=5mm of f0,font=\scriptsize] {$A_0,B_0\rightarrow S_0$};
    \node[below=5mm of f1,font=\scriptsize] {$A_1,B_1\rightarrow S_1$};
    \node[below=5mm of f2,font=\scriptsize] {$A_2,B_2\rightarrow S_2$};
    \node[below=5mm of f3,font=\scriptsize] {$A_3,B_3\rightarrow S_3$};
  \end{tikzpicture}
  \vspace{4mm}
\begingroup
\paragraph{リップルキャリー加算器}
        下位桁の$C_{out}$を次の桁の$C_{in}$へ渡します。
        単純ですが、桁数が増えると桁上がりの伝播遅延が増えます。

\par
\endgroup
\medskip
\begingroup
\paragraph{ALU}
        算術論理演算装置（ALU）は、加算、減算、論理演算、比較などを選択信号に応じて実行します。

\par
\endgroup
\par
\endgroup


\section{記憶回路}
\label{sec:01-digital-computer-10}

ALUは入力が来れば答えを出せますが、次の瞬間まで覚えてはいません。送信処理を何段階も進めるには、直前の結果や現在位置を残す必要があります。出力を入力へ戻すことで、回路に過去を持たせます。

\begingroup
\begingroup
\paragraph{組み合わせ回路の限界}
        出力が現在の入力だけで決まるため、入力が消えると過去の値を保持できません。

\paragraph{順序回路}
        出力を回路へ戻すフィードバックによって状態を保持します。
        出力は、現在の入力と直前までの状態で決まります。

\par
\endgroup
\medskip
\begingroup
      \centering
      \begin{tikzpicture}[node distance=9mm]
        \node[accentbox] (state) {状態記憶};
        \node[box,left=of state] (input) {入力};
        \node[box,right=of state] (output) {出力};
        \draw[flow] (input) -- (state);
        \draw[flow] (state) -- (output);
        \draw[softflow] (output.south) -- ++(0,-8mm) -| (state.south);
      \end{tikzpicture}
      \vspace{4mm}
\paragraph{フリップフロップ}
        SRフリップフロップはセットとリセットで状態を変えます。
        Dフリップフロップはクロックの有効なエッジで入力$D$を取り込み、次のエッジまで保持します。

\par
\endgroup
\par
\endgroup


\section{クロックとレジスタ}
\label{sec:01-digital-computer-11}

状態を保持できても、各部品がばらばらの時刻に書き換われば、途中の値を完成した結果と誤認します。クロックは更新の節目をそろえ、レジスタはその節目ごとにCPUがすぐ使う値を受け取ります。計算はここで時間の順序を持ちます。

\begingroup
\begingroup
\paragraph{クロック同期}
        組み合わせ回路が次の値を計算し、フリップフロップがクロックエッジで一斉に取り込みます。
        エッジ間には、信号が安定する時間が必要です。

\paragraph{レジスタ}
        複数のDフリップフロップを並べると、複数ビットを同時に保持するレジスタになります。

\par
\endgroup
\medskip
\begingroup
      \centering
      \begin{tikzpicture}[x=7mm,y=7mm]
        \draw[draw=DeepNavy] (0,0) -- (1,0) -- (1,1) -- (2,1) -- (2,0) -- (3,0) -- (3,1) -- (4,1) -- (4,0) -- (5,0);
        \foreach \x in {1,3}{\draw[draw=MutedBlue,dashed] (\x,-0.5) -- (\x,1.5);}
        \node[font=\scriptsize] at (2.5,-0.7) {クロック};
        \node[box,minimum width=20mm] at (2.5,-2) {状態更新};
      \end{tikzpicture}
      \vspace{3mm}
\paragraph{状態機械}
        レジスタが現在状態を保持し、組み合わせ回路が入力と現在状態から次状態を計算します。

\par
\endgroup
\par
\endgroup


\section{メモリ}
\label{sec:01-digital-computer-12}

レジスタは速い一方、会話文、注文データ、プログラム全体を置くには小さすぎます。多くの値へ番号を付け、必要な場所を選んで読み書きする仕組みがメモリです。入力された一文も、まずはアドレスを持つ領域へ収まります。

\begingroup
\begingroup
\paragraph{三つの信号群}
        メモリとのやり取りには、読み書きする位置を指定する\term{アドレス}、読み出す値または書き込む値を運ぶ\term{データ}、読み出しか書き込みかと動作時刻を指定する\term{制御信号}を使います。

\paragraph{バス}
        \term{バス}は、複数ビットの信号線をまとめた伝送路です。
        複数装置で共有する構成もあります。

\par
\endgroup
\medskip
\begingroup
      \centering
      \begin{tikzpicture}[node distance=1.5mm,every node/.append style={font=\small}]
        \node[accentbox,minimum width=27mm,minimum height=18mm] (mem) {メモリセル配列};
        \node[box,left=of mem] (dec) {行・列\\デコーダ};
        \node[box,above=of dec] (addr) {アドレス};
        \node[box,right=of mem] (data) {データ};
        \node[box,below=of dec] (ctrl) {制御};
        \draw[flow] (addr) -- (dec);
        \draw[flow] (dec) -- (mem);
        \draw[<->,draw=MutedBlue] (mem) -- (data);
        \draw[flow] (ctrl) -- (dec);
      \end{tikzpicture}
\par
\endgroup
  \source{MIT 6.111, “Memory Devices: SRAM / DRAM”}
\par
\endgroup


\section{SRAMとDRAM}
\label{sec:01-digital-computer-13}

すべての記憶を最速の回路で作れれば話は簡単ですが、速度、面積、消費電力、価格は両立しません。そこで、近くて速い記憶と、遠くて大きい記憶を別の方式で作ります。その代表がSRAMとDRAMです。

\begingroup
\begingroup
\paragraph{SRAM}
        代表的なSRAMセルは、トランジスタで構成したラッチに状態を保持します。
        電源がある間は定期的なリフレッシュを必要としません。

\paragraph{用途の傾向}
        高速なキャッシュなど、容量当たりの回路規模よりアクセス時間を重視する場所に使われます。

\par
\endgroup
\medskip
\begingroup
\paragraph{DRAM}
        セル内のコンデンサに蓄えた電荷で状態を表します。
        電荷が漏れるため、読み出しと再書き込みによるリフレッシュが必要です。

\paragraph{用途の傾向}
        高密度化しやすいため、主記憶に使われます。
        制御はSRAMより複雑です。

\par
\endgroup
  \takeaway{「SRAMは速い、DRAMは遅い」だけでなく、保持原理と密度の違いを押さえます。}
  \source{MIT 6.111, “Memory Devices: SRAM / DRAM”}
\par
\endgroup


\section{メモリ階層}
\label{sec:01-digital-computer-14}

異なる記憶をただ並べただけでは、プログラムが毎回置き場所を選ばなければなりません。実際には、最近使った値や近くの値を速い側へ保つことで、大きな記憶を持ちながら多くのアクセスを短時間で済ませます。抽象化の裏にある性能差が、初めてアプリの待ち時間へ顔を出します。

\begingroup
  \centering
  \begin{tikzpicture}[node distance=1mm,scale=0.82,transform shape]
    \node[accentbox,minimum width=35mm] (reg) {レジスタ：最小・最速};
    \node[box,minimum width=58mm,below=of reg] (l1) {L1キャッシュ};
    \node[box,minimum width=82mm,below=of l1] (l2) {L2／L3キャッシュ};
    \node[box,minimum width=106mm,below=of l2] (ram) {メインメモリ（DRAM）};
    \node[box,minimum width=130mm,below=of ram] (storage) {SSD／ストレージ：大容量・高遅延};
  \end{tikzpicture}
  \vspace{1mm}
\begingroup
\paragraph{局所性}
        直近に使ったデータや、その近くのデータを再び使う傾向を利用します。

\par
\endgroup
\medskip
\begingroup
\paragraph{キャッシュミス}
        上位層にデータがなければ、下位層からまとまり単位で取得します。
        この待ち時間が実行性能へ影響します。

\par
\endgroup
\par
\endgroup


\section{ノイマン型コンピュータ}
\label{sec:01-digital-computer-15}

計算する回路と記憶する回路がそろいました。次の一歩は、処理の手順そのものもデータと同じメモリへ置くことです。手順を交換可能なビット列にすると、同じ機械がブラウザにもAIアプリにもなれます。

\begingroup
\begingroup
      \centering
      \begin{tikzpicture}[node distance=5mm,every node/.append style={font=\small}]
        \node[accentbox,minimum width=48mm,minimum height=20mm] (cpu) {CPU\\制御装置・レジスタ・ALU};
        \node[box,right=of cpu,minimum height=20mm] (mem) {メモリ\\命令・データ};
        \node[box,below=of cpu] (io) {入出力装置};
        \draw[<->,draw=MutedBlue,line width=0.6pt] (cpu) -- node[above,font=\scriptsize] {アドレス・データ・制御} (mem);
        \draw[<->,draw=MutedBlue,line width=0.6pt] (cpu) -- (io);
      \end{tikzpicture}
\par
\endgroup
\medskip
\begingroup
\paragraph{プログラム内蔵方式}
        命令を数値としてメモリへ格納します。
        CPUは命令を順に読み出して処理します。

\paragraph{境界の意味}
        ハードウェアは命令セットの動作を実装し、ソフトウェアはその命令列で処理を記述します。

\par
\endgroup
\par
\endgroup


\section{CPUの命令実行}
\label{sec:01-digital-computer-16}

メモリに置かれた命令は、順番に読まれるだけでは仕事になりません。CPUは現在位置から命令を取り出し、そのビット列の意味を読み、必要な回路を選んで状態を更新します。送信ボタンの処理も、この短い周期の繰り返しとして進みます。

\begingroup
  \centering
  \begin{tikzpicture}[node distance=7mm]
    \node[accentbox] (fetch) {1　フェッチ};
    \node[box,right=of fetch] (decode) {2　デコード};
    \node[box,right=of decode] (exec) {3　実行};
    \node[box,right=of exec] (write) {4　書き戻し};
    \draw[flow] (fetch) -- (decode);
    \draw[flow] (decode) -- (exec);
    \draw[flow] (exec) -- (write);
    \draw[softflow] (write.south) -- ++(0,-8mm) -| node[below,pos=0.25,font=\scriptsize] {PCを更新} (fetch.south);
  \end{tikzpicture}
  \vspace{6mm}
\begingroup
\paragraph{フェッチとデコード}
        プログラムカウンタ（PC）が示すアドレスから命令を読みます。
        制御装置は命令の演算種別とオペランドを解釈します。

\par
\endgroup
\medskip
\begingroup
\paragraph{実行と書き戻し}
        ALU、レジスタ、メモリ、分岐回路を動かし、結果を保存します。
        実際のCPUはパイプラインなどで複数命令を重ねて処理できます。

\par
\endgroup
  \source{東京大学「コンピュータハードウェア」第7回}
\par
\endgroup


\section{加算命令の実行}
\label{sec:01-digital-computer-17}

ここでCPUの一周期を止め、加算命令一つを拡大します。レジスタから値が出て、ALUへ入り、結果が別のレジスタへ戻るまでを追えば、「命令を実行する」という言葉を回路上の移動へ結び付けられます。

\begingroup
\noindent{\sffamily\bfseries\color{MutedBlue}例：\code{ADD R3, R1, R2} はR1とR2を加え、R3へ保存します}\par\medskip
\begingroup
      \centering
      \begin{tikzpicture}[node distance=5mm]
        \node[box] (r1) {R1};
        \node[box,below=of r1] (r2) {R2};
        \node[accentbox,right=14mm of $(r1)!0.5!(r2)$] (alu) {ALU\\加算};
        \node[box,right=of alu] (r3) {R3};
        \draw[flow] (r1.east) -- (alu.west);
        \draw[flow] (r2.east) -- (alu.west);
        \draw[flow] (alu) -- (r3);
        \node[box,above=of alu] (ctrl) {制御装置};
        \draw[softflow] (ctrl) -- node[right,font=\scriptsize] {演算選択} (alu);
      \end{tikzpicture}
\par
\endgroup
\medskip
\begingroup
      加算命令は、次の順序でデータ経路を使います。
      \begin{enumerate}
        \item 命令をメモリから読みます。
        \item 命令形式からR1、R2、R3を特定します。
        \item レジスタR1とR2を読みます。
        \item ALUで加算します。
        \item 結果をR3へ書き戻します。
        \item PCを次の命令へ進めます。
      \end{enumerate}
\par
\endgroup
  \takeaway{命令は、データ移動と回路選択を同期させた処理です。}
\par
\endgroup


\section{命令の先へ}
\label{sec:01-digital-computer-18}

連続的な電圧はしきい値によって0と1になり、ビット列は規則によって数や文字になりました。トランジスタから組み上げたゲートは計算と記憶を担い、CPUはメモリ中の命令を読みながら状態を更新します。送信という上位の出来事も、最下層ではこの物理的な変化の連続です。

ただし、ここには一つ大きな空白があります。CPUが読んだのは機械向けの命令ですが、開発者が書いたのは関数や型を持つコードでした。次章では視点を少し上げ、人間の意図がどのように命令の列へ翻訳されたのかをたどります。

\takeaway{物理から計算への流れは、電圧を論理値として読み、論理値を表現と回路へ組み、回路の状態変化を命令として制御する積み重ねです。}


\section{旅をたどり直す}
\label{sec:01-digital-computer-19}

送信ボタンの操作がCPUの状態変化になるまでを思い浮かべ、次の問いへ答えてください。用語を一つずつ定義するだけでなく、前後の仕組みへどう渡るかを説明します。

\begingroup
  \begin{enumerate}
    \item Highを「5 Vそのもの」と説明すると不正確なのはなぜですか。
    \item $1011\,0110_2$を16進数へ変換してください。
    \item 8ビットの2の補数で表せる範囲と、$-18$のビット列を示してください。
    \item NANDだけでORを構成する式を示してください。
    \item 組み合わせ回路と順序回路の出力を決める情報は何ですか。
    \item SRAMとDRAMは、状態の保持方法がどのように異なりますか。
    \item \code{ADD R3, R1, R2} の実行で、レジスタとALUはどの順に使われますか。
  \end{enumerate}
\par
\endgroup


\section{手を動かす：数の表現}
\label{sec:01-digital-computer-20}

同じビット列でも、幅と解釈規則が変われば見える値は変わります。短いプログラムを観測窓にして、回路の表現がアプリケーションの計算結果へ現れる瞬間を確かめます。

\begingroup
\begingroup
      \begin{lstlisting}[language=Python]
def u8(x):
    return x & 0xff

for x in [127, 128, 255, 256, -1, -5]:
    print(x, format(u8(x), "08b"), u8(x))

print(0.1 + 0.2)
print((0.1 + 0.2) == 0.3)
      \end{lstlisting}
\par
\endgroup
\medskip
\begingroup
\paragraph{観察する項目}
        実行結果では、次の四点を観察します。
        \begin{itemize}
          \item 256の下位8ビット
          \item $-1$と$-5$のビット列
          \item 0.1と0.2の加算結果
          \item 厳密な等値比較の結果
        \end{itemize}

\paragraph{説明課題}
        出力を記録し、各結果を「固定幅」と「2進表現」の二語を使って説明してください。

\par
\endgroup
\par
\endgroup
